\documentclass[11pt]{article}

%%METADATA
\title{GA1025 Macro Theory \\Solution to Problem Set 9}
\author{
Junbiao Chen\thanks{E-mail: jc14076@nyu.edu.}
}
\date{\today}


%%PACKAGES
\usepackage{mdframed} % For boxed environments
\usepackage{graphicx}
\usepackage{grffile}
\usepackage{tabularx}
\usepackage{setspace}
\usepackage{amsmath,amsthm,amssymb}
\usepackage[hyphens]{url}
\usepackage{natbib}
\usepackage[font=normalsize,labelfont=bf]{caption}
\usepackage[margin=1in]{geometry}
\usepackage{hyperref}
\hypersetup{colorlinks=true,urlcolor=blue,citecolor=blue}
\usepackage{stmaryrd}  %Package with \boxast command
\usepackage{enumerate}% http://ctan.org/pkg/enumerate %Supports lowercase Roman-letter enumeration
\usepackage{verbatim} %Package with \begin{comment} environment
%\usepackage{enumitem}
\usepackage{physics}
\usepackage{tikz}
\usepackage{listings}
\usepackage{upquote}
\usepackage{booktabs} %Package with \toprule and \bottomrule
\usepackage{etoc}     %Package with \localtableofcontents
\usepackage{placeins}    %Package that prevent repositioning the tables
\usepackage{multicol}
\usepackage{bm}
\usepackage{subfig}
\usepackage{csquotes}

\definecolor{dkgreen}{rgb}{0,0.6,0}
\definecolor{gray}{rgb}{0.5,0.5,0.5}
\definecolor{mauve}{rgb}{0.58,0,0.82}

\lstset{language=bash,
  frame=tb,
  aboveskip=3mm,
  belowskip=3mm,
  showstringspaces=false,
  columns=flexible,
  basicstyle={\small\ttfamily},
  numbers=none,
  numberstyle=\tiny\color{gray},
  keywordstyle=\color{blue},
  commentstyle=\color{dkgreen},
  stringstyle=\color{mauve},
  breaklines=true,
  breakatwhitespace=false,
  tabsize=3
}

\lstset{language=C,
  aboveskip=3mm,
  belowskip=3mm,
  showstringspaces=false,
  columns=flexible,
  basicstyle={\small\ttfamily},
  numbers=none,
  numberstyle=\tiny\color{gray},
  keywordstyle=\color{blue},
  commentstyle=\color{dkgreen},
  stringstyle=\color{mauve},
  breaklines=true,
  breakatwhitespace=false,
  tabsize=4
}

\definecolor{lightblue}{rgb}{0.68, 0.85, 0.9} %

%CUSTOM DEFINITIONS
\theoremstyle{definition}
\newtheorem{definition}{Definition}[section]
\newtheorem*{remark}{Remark}
\setcounter{secnumdepth}{3}
\usepackage{tikz}
\usetikzlibrary{arrows.meta}
\usetikzlibrary{automata, positioning, arrows, calc}

\tikzset{
	->,  % makes the edges directed
	>=stealth, % makes the arrow heads bold
	shorten >=2pt, shorten <=2pt, % shorten the arrow
	node distance=3cm, % specifies the minimum distance between two nodes. Change if n
	every state/.style={draw=blue!55,very thick,fill=blue!20}, % sets the properties for each ’state’ n
	initial text=$ $, % sets the text that appears on the start arrow
 }

%% PROPOSITION
% Define the Proposition environment
\newmdenv[
  innerleftmargin=10pt, 
  innerrightmargin=10pt,
  innertopmargin=10pt,
  innerbottommargin=10pt,
  linecolor=black, 
  linewidth=1pt,
  backgroundcolor=white, 
  roundcorner=5pt
]{propositionbox}

\newtheoremstyle{boldtitle} % Define a new theorem style
  {10pt} % Space above
  {10pt} % Space below
  {\itshape} % Body font
  {} % Indent amount
  {\bfseries} % Theorem head font
  {.} % Punctuation after theorem head
  { } % Space after theorem head
  {} % Theorem head spec

\theoremstyle{boldtitle} % Use the custom style
\newtheorem{proposition}{Proposition} % Define the proposition environment

% Redefine the proposition environment to use the box
\newenvironment{boxedproposition}[1][]
{\begin{propositionbox}\begin{proposition}[#1]}
{\end{proposition}\end{propositionbox}}

%%FORMATTING
\usepackage[bottom]{footmisc}
\onehalfspacing
\numberwithin{equation}{section}
\numberwithin{figure}{section}
\numberwithin{table}{section}
\bibliographystyle{../bib/aeanobold-oxford}


%main text
\begin{document}

\maketitle


\section{Exercise 9.7 Two agents}
\subsection{a. Answer} 
Denote the transition matrix $P$, we have 
\[
\pi_{t+1}' = \pi_t' P
\]
It follows that 
\[
\pi_t' = \pi_0' P^t.
\]

\subsection{b. Answer} 
The Largrangian is given by 
\[
\mathcal{L} =  \sum_{t=0}^\infty \sum_{s^t} \beta^t  pi_t(s^t)  \bigg(\theta \ln \left[ c_t^1 (s^t) \right] + (1 - \theta) \ln \left[ c_t^2 (s^t) \right] \bigg) 
+ \mu \sum_{t=0}^\infty \sum_{s^t} \bigg(Y_t(s_t) - c_t^1 (s^t) - c_t^2 (s^t) \bigg)
\]
The FOCs are:
\begin{align*}
    [c^1_t(s^t)]: \quad \mu & = \frac{\theta \beta^t \pi_t(s^t) }{c_t^1(s^t)} \\
    [c^2_t(s^t)]: \quad \mu & = \frac{(1 - \theta) \beta^t \pi_t(s^t) }{c_t^2(s^t)}  \\
\end{align*}
Applysing the market clearing condition: $Y_t = 1 + s_t = c_t^1(s^t) + c_t^2(s^t)$, we have 
\[
\frac{\theta  }{c_t^1(s^t)} = \frac{1 - \theta  }{1+ s_t - c_t^1(s^t)} 
\]
Therefore, the solution is given by 
\begin{align*}
    c^1_t(s^t) & =  \theta (1 + s_t) \\
    c^2_t(s^t) & =  (1 - \theta) (1 + s_t) \\
\end{align*}


\subsection{c. Answer} 
The competitive equilibrium with time-0 trading is given by 
an allocation $\{c_t^1(s^t), c_t^2(s^t)\}$ and a price system $q_t^0(s^t)$, such that 
\begin{enumerate}
    \item the allocation solves each consumer's problem,
    \item market clears: ${c_t(s^t)}^1 + {c_t(s^t)}^2 = {y_t(s^t)}^1 + {y_t(s^t)}^2 = Y_t, \quad \forall t, s^t$. 
\end{enumerate}


\subsection{d. Answer} 
The FOCs for agents 1 and 2 imply that
\begin{align*}
    \beta^t \pi_t(s^t) [c_t^1(s^t)]^{-1} &= \mu^1 q_t^0(s^t) \\
    [c_0^1(s_0)]^{-1} &= \mu^1 
\end{align*}

\begin{align*}
    \beta^t \pi_t(s^t) [c_t^2(s^t)]^{-1} &= \mu^2 q_t^0(s^t) \\
    [c_0^2(s_0)]^{-1} &= \mu^2
\end{align*}
where I noramlize $q_0^0(s_0) = 1$ and use the facts that $\beta^0 =1$, and $\pi_0(s_0)=1$. 

Adding these equations together, we have 
\[
\beta^t \pi_t(s^t)Y_0(s_0) = Y_t(s_t) q_t^0(s^t)
\]
Therefore, we have
\[
q_t^0(s^t) = \frac{\beta^t \pi_t(s^t)Y_0(s_0)}{Y_t(s_t) } = \frac{\beta^t \pi_t(s^t)(1+s_0)}{1 + s_t} 
\]


\subsection{e. Answer} 
From (d), we have 
\[
\beta^t \pi_t(s^t) [c_t^1(s^t)]^{-1} = \mu^1 \frac{\beta^t \pi_t(s^t)(1+s_0)}{1 + s_t} 
\]
Massaging terms on both sides, we have 
\[
c_t^1(s^t) = \frac{1}{\mu^1 (1 + s_0)} (1 + s_t),
\]
which coincides with the Pareto solution if $\theta = \frac{1}{\mu^1 (1 + s_0)}$.

\noindent \textbf{Remark:} The First Welfare Theorem implies that C.E. is Pareto optimal if there is no externalities.


\subsection{f. Answer} 
In the competitive equilibrium, a price system is determined by inidividuals' MRS (up to some normalization and proper discounting.)


\subsection{g. (Recursive competitive equlibrium) Answer} 
Let $v^i(a,s)$ be the optimal value of agent $i$ with states $a$ (wealth) and $s$ (aggregate state).
The optimal value function satisfies the Bellman equation
\begin{align*}
v^i(a, s) & = \max_{s',a'(s')}\quad \ln c^i(s) + \beta \sum_{s'} \pi(s'|s) v^i(a'(s'),s') \\
\text{subject to } \quad & c^i(s) + \sum_{s'} Q(s'|s) a'(s') \leq y^i(s) + a \\ 
& -a'(s') \leq A^i (s'), \forall s' \in S
\end{align*}
where $A^i(s')$ is state-by-state debt limits given by the following recursive:
\[
A^i(s) = y^i(s) + \sum_{s'} Q(s'|s)  A^i(s'), \quad i = 1, 2
\]

\noindent \textbf{Definition}
A recursive competitive equilibrium is given by 
an initial distribution of wealth $\{a_0^i\}$, a set of borrowing limits 
$A^i(s)$, a pricing kernel $Q(s'|s)$, 
a set of value functions $\{v^i(a, s) \}$, and 
decision rules $\{c^i(s), {a'}^i(s) \}$ such that 
\begin{enumerate}
    \item The state-by-state borrowing constraints satisfy the recursion: $  A^i(s) = y^i(s) + \sum_{s'} Q(s'|s)  A^i(s')$.
    \item For all $i = 1, 2$, given $a^i_0$, $A^i(s)$, and $Q(s'|s)$, the value function and decision rules solve the consumers' problems.
    \item For all realizations of $\{s_t \}_{t=0}^\infty$, the consumption and asset portfolios $\{ \{c^i_t, \{{{a_{t+1}'}^i}(s')\}_{s'} \}_i\}_t$ 
    satisfy 
    \begin{align*}
        c^1_t + c^2_t & = y^1(s_t) + y^2(s_t) \\
        {a'_{t+1}}^1 + {a'_{t+1}}^2 & = 0
    \end{align*}
    for all $t$ and $s'$.
    \item The initial wealth distribution satisfies $a_0^1 + a_0^2 = 0$.
\end{enumerate}


\subsection{h. Answer} 
Using the FOCs of the problem and a Benveniste-Scheinkman formula, we have 
\[
Q(s_{t+1}|s_t) = \frac{\beta [c^i_{t+1}(s_{t+1})]^{-1} \pi(s_{t+1} | s_t)
    }{ [c^i(s_t)]^{-1}}
\]
Given the parameterization of $\beta = 0.95$ and 
$P = \begin{bmatrix}
    0.8 & 0.2 \\ 0.3 & 0.7
\end{bmatrix} $, we have 

\begin{table}[h!]
\centering
\begin{tabular}{ccc}
\hline
$s_t$ & $s_{t+1}$ & $Q(s_{t+1} | s_t)$ \\ \hline
0     & 0         & $0.95 \cdot 0.8$ \\ 
0     & 1         & $0.95 \cdot \frac{c^1(0)}{c^1(1)} \cdot 0.2$ \\  
1     & 0         & $0.95 \cdot \frac{c^1(1)}{c^1(0)} \cdot 0.3$ \\  
1     & 1         & $0.95 \cdot 0.7$ \\ \hline
\end{tabular}
\caption{Transition probabilities and corresponding Q-values.}
\label{tab:q_values}
\end{table}

Therefore, using No-Arbitrage condition, 
the prices of one-period Arrow securities are given by 
\[
p(s) = \sum_{s_{t+1}} Q(s_{t+1}|s_t = s)
\]

%% =========== %% ========= %%
\bibliography{../bib/references.bib}

\end{document}